%% introduction.tex
%%

%% ==============================
\section{Introduction}
\label{sec:Introduction}
%% ==============================

Kubernetes' built-in autoscalers -- the Horizontal Pod Autoscaler (HPA) and the event-driven KEDA -- are reactive: they scale only once a resource metric has already crossed a threshold, which under bursty, cyclical workloads causes either SLA-violating under-provisioning during the delay or costly over-provisioning afterward. Proactive, ML-based forecasting is a natural fix, but it introduces a new risk: acting on every noisy prediction can make the system more agile at the cost of being less stable.

Wanigasooriya \& Ekanayake \cite{Wanigasooriya26} demonstrate exactly this with NimbusGuard, a Deep Q-Network + LSTM proactive autoscaler evaluated against reactive HPA and KEDA on a real Kubernetes testbed. NimbusGuard wins on responsiveness, but their own results table shows it running with the highest average replica count and double the total scaling events of either reactive baseline -- described in their own words as ``the most agile and least stable system.'' They name addressing this stability cost as future work.

\textbf{Research question:} can a proactive autoscaler keep most of its SLA-violation advantage over reactive HPA while substantially reducing the instability (scaling-event frequency, pod-count volatility) that an unfiltered, DQN-style proactive policy exhibits?

\textbf{Objectives:}
\begin{itemize}
  \item Reproduce the trade-off NimbusGuard reports -- proactive scaling beats reactive HPA on SLA compliance, at a cost in over-provisioning and instability -- with a simpler, fully reproducible feed-forward predictor rather than a full DQN+LSTM+LLM pipeline.
  \item Design and build the stability fix NimbusGuard names but does not build: smoothing and hysteresis applied to the same underlying forecast.
  \item Quantify the trade-off honestly across multiple metrics and multiple random seeds, rather than reporting only the metrics that favour the new design.
\end{itemize}

The remainder of this report is structured as follows: Section~\ref{sec:rw} reviews the baseline paper and the gap it names; Section~\ref{sec:Design} describes the three scaling policies compared; Section~\ref{sec:Implementation} describes the workload simulator, training pipeline, and deployment; Section~\ref{sec:Evaluation} reports results across 5 seeds; Section~\ref{sec:Conclusion} discusses the trade-off and future work.
